\section{10 June 2026}

Meeting assets for mtg: Dean PASS are ready! 

\textbf{Meeting summary }

\subsection{Quick recap}

This meeting was a PhD supervision discussion between Professor Fares, Dr. Hozefa (Dr. Jizwada), and DeanPretorius regarding the development of a PhD research plan focused on space robotics and locomotion. The conversation covered DeanPretorius's background in systems engineering and industry experience, including work on mining robotics and IoT projects at CSIR in South Africa. Professor Fares shared his recent contact with a space robotics specialist from DLR Germany and discussed plans to potentially invite him for a talk in October or November. The group explored several technical directions including differential geometry, Koopman operators for linearizing nonlinear systems, and friction modeling in robotic locomotion. DeanPretorius proposed integrating systems engineering approaches using NASA's Systems Engineering Handbook with real-world robot applications, including edge-to-cloud computing and virtual reality components. Dr. Hozefa provided feedback that the proposed research direction was feasible and offered insights about Koopman operators and their applications in fluid mechanics and robotics. The discussion also touched on the potential for co-supervision arrangements and the importance of creating compelling demonstration videos to showcase research findings.

\subsection{Next steps}

\subsubsection{DeanPretorius}

\begin{itemize}
    \item Create an Overleaf document titled "PSD Dean's PSD Plan" with separate LaTeX pages for tentative plan, ideas, and other relevant sections; document thoughts, comments, and progress discussed in meetings.
    \item Summarize today's discussion points and add them to the Overleaf document.
    \item Investigate the feasibility and novelty of applying Koopman operators (and related methods such as data-driven Koopman/EDMD) to linearize or analyze the nonlinear dynamics in the proposed robotic system, and check the current state of the art (including flow matching paper shared by Hozefa).
    \item Review the flow matching paper shared by Hozefa and consider its relevance to the PhD plan.
    \item Continue thinking about and developing the PhD plan, particularly on locomotion (and optionally manipulation), and divide time among the suggested priorities.
    \item Use NYU email/chat for future communications and document sharing with Fares.
    \item When logistics are clear, coordinate with Fares regarding lab space and desk arrangements.
    \item If considering co-supervision, first identify clear limitations or needs before proposing specific co-supervisors, and discuss with Fares if/when appropriate.
\end{itemize}

\subsubsection{Fares}

\begin{itemize}
    \item Complete exploration of differential geometry (including Grassmannian and other approaches) and update findings in the shared Overleaf document or present in the next meeting.
    \item Send Zoom meeting minutes to DeanPretorius after the meeting.
\end{itemize}

\subsection{Summary}

\subsubsection{Academic Updates and Career Discussions}

The meeting began with congratulations to Hozefa for recently completing his PhD, which he described as feeling more like relief than happiness after the struggle of the process. DeanPretorius shared his work at CSIR in South Africa, where he does research for government applications, and mentioned that his role is similar to an HOD with a Palestinian colleague. Fares expressed interest in CSIR, noting it as a prestigious African institution, and shared his own consideration of moving to South Africa for a career change.

\subsubsection{Space Robotics Collaboration Planning}

Fares reported creating a new contact with a space robotics specialist from DLR Germany and plans to invite him for a talk between October and December next semester. Fares suggested the possibility of creating a visit for DeanPretorius to the German research center as part of potential collaboration. They discussed the importance of applying theoretical work to industrial scenarios to make the PhD more industry-relevant, with Fares emphasizing the need to create their own industrial applications regardless of external collaboration.

\subsubsection{Robot Collaboration and Testing Planning}

Fares and DeanPretorius discussed potential collaboration opportunities, including using existing facilities like Khalifa University's ARENA program or creating their own testing arena for robot locomotion research. They agreed that simulation would be essential but not sufficient, and that real system testing would be necessary to validate their work. Fares assigned DeanPretorius to create a LaTeX-based project management document on Overleaf to better track progress and discussions, and DeanPretorius confirmed his familiarity with LaTeX. The conversation also briefly touched on DeanPretorius's previous industry experience, including a mining safety project involving a handheld rock detection device.

\subsubsection{Systems Engineering PhD Research Discussion}

DeanPretorius discussed his background, including his work at the CSIR where they conduct research and publish while avoiding competition with universities, taking projects from TRL 3-4 to 9. He explained his interest in applying systems engineering principles to his PhD research, specifically wanting to use the NASA Systems Engineering Handbook as a framework for requirements gathering, use case analysis, and stakeholder management. DeanPretorius shared an example from his industry experience where a missing power button in a rock detector highlighted the importance of considering user requirements beyond laboratory testing.

\subsubsection{PhD Research Directions Discussion}

The group discussed research directions for DeanPretorius's PhD, focusing on locomotion in space and sensorless friction modeling. Fares suggested exploring differential geometry and Koopman operators for linearizing nonlinear systems, while DeanPretorius proposed using edge-to-cloud computing with a Spot robot to model friction dynamics in different environments. Hozefa confirmed the feasibility of the proposed PhD direction and advised that research often requires pivoting based on findings, suggesting DeanPretorius should document current ideas and explore Koopman operators further. The discussion also covered practical aspects of co-supervision, with Hozefa cautioning that additional supervisors should complement rather than conflict with the main supervisor's direction.

\section{ 22 May 2026}

\textbf{Meeting summary }

\subsection{Quick recap}

The meeting focused on discussing DeanPretorius's research on applying Grassmannian manifolds to robotic friction modeling. DeanPretorius presented his previous work on calculating friction in 3D space using optimization methods, which he had implemented on monopod and biped robots. The team discussed how to extend this work to multi-legged robots like hexapods or octopods using Grassmannian manifolds to reduce computational complexity. Hozefa suggested that friction cones themselves are already Riemannian manifolds, providing a potential connection to the Grassmannian approach. The group agreed that DeanPretorius should research the state of the art in applying Grassmannian manifolds to friction modeling and friction cones, and share his thesis and ICRA paper with the team for further analysis.

\subsection{Next steps}

\subsubsection{DeanPretorius}

\begin{itemize}
    \item Search the literature/state of the art for existing work on the application of Grassmannian or Riemannian manifolds to friction cones, sliding, or related topics in robotics.
    \item Use AI tools (e.g., NotebookLM, ChatGPT, Claude) to explore research gaps and possible approaches for applying Grassmannian manifold concepts to the friction optimization problem.
    \item Send the link to his ICRA paper and thesis via email to Fares (and optionally Hozefa).
\end{itemize}

\subsubsection{Fares}

Review the ICRA paper and thesis sent by DeanPretorius before the next meeting.

\subsubsection{Collaboration}

All: Schedule and attend the next weekly meeting (noted to be half an hour long from next week).

\subsection{Summary}

\subsubsection{DeanPretorius PhD Program Introduction}

Fares introduced DeanPretorius, who is joining NYU's PhD program in the space systems engineering track, and mentioned that Dean had discussed his ICRA paper during the hiring process. Fares suggested there was potential to extend Dean's work using Grassmann and Manifold methods, and proposed involving Hozefa, a postdoc in Fares' lab with experience in Grassmann Manifold, to assist with this work.

\subsubsection{Admission and Teaching Arrangements Update}

DeanPretorius confirmed that his clearance and transcript are complete, and he is waiting for the work permit to be issued, which should allow the university to start the admission process. The expected arrival date is in mid-August, coinciding with the start of the semester. Fares mentioned he will be teaching a space robotics course next semester and suggested continuing their meetings weekly with a reduced duration of 30 minutes. DeanPretorius reported progress on reading assigned papers and mentioned technical difficulties with his office Wi-Fi and microphone, which were later resolved.

\subsubsection{Grassmannian Mathematics Project Preparation}

DeanPretorius discussed his preparation for a project involving Grassmannian mathematics, including reading relevant papers and creating summaries using NotebookLM. He explained his understanding of the Grassmannian and its applications, as well as his previous master's work on calculating friction in 3D space. DeanPretorius mentioned having access to the necessary papers and thesis but faced technical difficulties accessing Google Drive due to potential firewall issues.

\subsubsection{Multi-Legged Robot Expansion Discussion}

Dean Pretorius discussed expanding the Grismanian method from a single-legged robot to a multi-legged robot, initially working on a biped. He explained that the optimization method was chosen because friction performance during transient phases of motion is highly non-linear and difficult to compute. The next phase involves potentially working with an octopod or hexapod robot, using the Grismanian manifold to calculate net friction on multi-legged robots.

\subsubsection{Robot Motion Optimization Research}

DeanPretorius explained his research on optimizing robot motion by modeling friction constraints using the maximum dissipation principle, which allows for more accurate cone modeling compared to the traditional pyramid approach. He demonstrated the method's feasibility on monopod and biped robots, showing a reduction in computational power. The discussion concluded with plans to extend this work to hexapod or octopod robots, with Dr. Obidaka suggesting the use of Grismannian methods to further reduce the problem's complexity.

\subsubsection{Biped Simulation Force Optimization}

DeanPretorius explained a symbol representing ground contact force parallel to motion direction in a biped simulation. He discussed challenges with the non-linear optimization problem, including long simulation times (18,000 epochs taking six hours each) and issues with local minima. The conversation ended with Hozefa asking about the current solver being used for the optimization problem, but this part was cut off in the transcript.

\subsubsection{Trajectory Optimization Approach Discussion}

DeanPretorius explained his approach to trajectory optimization, which involves solving for the full lifecycle of motion using an epsilon relaxation method. He demonstrated how the solution space is iteratively tightened, starting with a large error margin (blue line) and gradually reducing it to find feasible solutions. The discussion included analysis of friction cone constraints and slip dynamics, with particular attention to how these factors affect motion both on Earth and in space. DeanPretorius clarified that this work represents classical trajectory optimization rather than model predictive control, as it optimizes the entire trajectory from rest to steady state.

\subsubsection{Robot Velocity Tracking Research}

DeanPretorius presented work on tracking robot forward velocity, capturing acceleration trajectories, and identifying heuristics in transient motion using a sliding mass model. He explained how they plotted the velocity of the center of gravity over time and noted similarities to a normal sliding box on a frictional surface. Fares then discussed the concept of considering contact points as projections on a Grassmannian manifold rather than separate spaces.

\subsubsection{Grassmannian Spaces for Robot Configurations}

The team discussed using Grassmannian spaces to generalize robot configurations and projections. Hozefa explained that each joint configuration and phase of motion (friction, sliding, and stopping) would create different manifolds that could be projected as points on the Grassmannian. Fares clarified that friction itself represents a projection, and DeanPretorius confirmed his understanding that the Grassmannian manifold represents the solution space and its curvature.

\subsubsection{Riemannian Manifolds in Optimization}

Fares explained the concept of Riemannian manifolds and their relevance to modeling solution spaces in optimization problems, particularly in the context of friction projection. DeanPretorius confirmed understanding and agreed to research existing literature on friction projection and Riemannian manifolds. Fares emphasized the exploratory nature of their work, noting the lack of clear solutions in the literature and the need for further research.

\subsubsection{AI Tools for Academic Work}

The participants discussed various AI tools for academic work, with DeanPretorius demonstrating NotebookLM's ability to create audio overviews from articles and Dr. Abudakar's use of it for student onboarding. Fares shared his experience using multiple AI tools including Claude, Gemini, ChatGPT, and DeepSeek, noting that while DeepSeek had superior search capabilities for niche topics, server reliability issues led him to switch to other platforms. The conversation ended with DeanPretorius attempting to demonstrate something that wasn't working properly at the time.

\subsubsection{Physics Simulations Research Discussion}

The team discussed DeanPretorius's work on physics simulations in Google Colab, focusing on friction cones and Riemannian manifolds. Hozefa suggested exploring how multi-contact scenarios work with individual cones and manifolds for each contact point. DeanPretorius agreed to research the application of Christmanian manifolds in friction space and identify research gaps. The team agreed to review DeanPretorius's thesis and paper, with Fares requesting the documents be sent by email for further analysis before the next meeting.

 
\section{13 May 2026}

\textbf{Meeting summary }

\subsection{Quick recap}

The meeting focused on discussing DeanPretorius's research on applying Grassmannian manifolds to robotic friction modeling. DeanPretorius presented his previous work on calculating friction in 3D space using optimization methods, which he had implemented on monopod and biped robots. The team discussed how to extend this work to multi-legged robots like hexapods or octopods using Grassmannian manifolds to reduce computational complexity. Hozefa suggested that friction cones themselves are already Riemannian manifolds, providing a potential connection to the Grassmannian approach. The group agreed that DeanPretorius should research the state of the art in applying Grassmannian manifolds to friction modeling and friction cones, and share his thesis and ICRA paper with the team for further analysis.

\subsection{Next steps}

\subsubsection{DeanPretorius}

\begin{itemize}
    \item Search the literature/state of the art for existing work on the application of Grassmannian or Riemannian manifolds to friction cones, sliding, or related topics in robotics.
    \item Use AI tools (e.g., NotebookLM, ChatGPT, Claude) to explore research gaps and possible approaches for applying Grassmannian manifold concepts to the friction optimization problem.
    \item Send the link to his ICRA paper and thesis via email to Fares (and optionally Hozefa).
\end{itemize}

\subsubsection{Fares}

Review the ICRA paper and thesis sent by DeanPretorius before the next meeting.

\subsubsection{Collaboration}

All: Schedule and attend the next weekly meeting (noted to be half an hour long from next week).

\subsection{Summary}

\subsubsection{DeanPretorius PhD Program Introduction}

Fares introduced DeanPretorius, who is joining NYU's PhD program in the space systems engineering track, and mentioned that Dean had discussed his ICRA paper during the hiring process. Fares suggested there was potential to extend Dean's work using Grassmann and Manifold methods, and proposed involving Hozefa, a postdoc in Fares' lab with experience in Grassmann Manifold, to assist with this work.

\subsubsection{Admission and Teaching Arrangements Update}

DeanPretorius confirmed that his clearance and transcript are complete, and he is waiting for the work permit to be issued, which should allow the university to start the admission process. The expected arrival date is in mid-August, coinciding with the start of the semester. Fares mentioned he will be teaching a space robotics course next semester and suggested continuing their meetings weekly with a reduced duration of 30 minutes. DeanPretorius reported progress on reading assigned papers and mentioned technical difficulties with his office Wi-Fi and microphone, which were later resolved.

\subsubsection{Grassmannian Mathematics Project Preparation}

DeanPretorius discussed his preparation for a project involving Grassmannian mathematics, including reading relevant papers and creating summaries using NotebookLM. He explained his understanding of the Grassmannian and its applications, as well as his previous master's work on calculating friction in 3D space. DeanPretorius mentioned having access to the necessary papers and thesis but faced technical difficulties accessing Google Drive due to potential firewall issues.

\subsubsection{Multi-Legged Robot Expansion Discussion}

Dean Pretorius discussed expanding the Grismanian method from a single-legged robot to a multi-legged robot, initially working on a biped. He explained that the optimization method was chosen because friction performance during transient phases of motion is highly non-linear and difficult to compute. The next phase involves potentially working with an octopod or hexapod robot, using the Grismanian manifold to calculate net friction on multi-legged robots.

\subsubsection{Robot Motion Optimization Research}

DeanPretorius explained his research on optimizing robot motion by modeling friction constraints using the maximum dissipation principle, which allows for more accurate cone modeling compared to the traditional pyramid approach. He demonstrated the method's feasibility on monopod and biped robots, showing a reduction in computational power. The discussion concluded with plans to extend this work to hexapod or octopod robots, with Dr. Obidaka suggesting the use of Grismannian methods to further reduce the problem's complexity.

\subsubsection{Biped Simulation Force Optimization}

DeanPretorius explained a symbol representing ground contact force parallel to motion direction in a biped simulation. He discussed challenges with the non-linear optimization problem, including long simulation times (18,000 epochs taking six hours each) and issues with local minima. The conversation ended with Hozefa asking about the current solver being used for the optimization problem, but this part was cut off in the transcript.

\subsubsection{Trajectory Optimization Approach Discussion}

DeanPretorius explained his approach to trajectory optimization, which involves solving for the full lifecycle of motion using an epsilon relaxation method. He demonstrated how the solution space is iteratively tightened, starting with a large error margin (blue line) and gradually reducing it to find feasible solutions. The discussion included analysis of friction cone constraints and slip dynamics, with particular attention to how these factors affect motion both on Earth and in space. DeanPretorius clarified that this work represents classical trajectory optimization rather than model predictive control, as it optimizes the entire trajectory from rest to steady state.

\subsubsection{Robot Velocity Tracking Research}

DeanPretorius presented work on tracking robot forward velocity, capturing acceleration trajectories, and identifying heuristics in transient motion using a sliding mass model. He explained how they plotted the velocity of the center of gravity over time and noted similarities to a normal sliding box on a frictional surface. Fares then discussed the concept of considering contact points as projections on a Grassmannian manifold rather than separate spaces.

\subsubsection{Grassmannian Spaces for Robot Configurations}

The team discussed using Grassmannian spaces to generalize robot configurations and projections. Hozefa explained that each joint configuration and phase of motion (friction, sliding, and stopping) would create different manifolds that could be projected as points on the Grassmannian. Fares clarified that friction itself represents a projection, and DeanPretorius confirmed his understanding that the Grassmannian manifold represents the solution space and its curvature.

\subsubsection{Riemannian Manifolds in Optimization}

Fares explained the concept of Riemannian manifolds and their relevance to modeling solution spaces in optimization problems, particularly in the context of friction projection. DeanPretorius confirmed understanding and agreed to research existing literature on friction projection and Riemannian manifolds. Fares emphasized the exploratory nature of their work, noting the lack of clear solutions in the literature and the need for further research.

\subsubsection{AI Tools for Academic Work}

The participants discussed various AI tools for academic work, with DeanPretorius demonstrating NotebookLM's ability to create audio overviews from articles and Dr. Abudakar's use of it for student onboarding. Fares shared his experience using multiple AI tools including Claude, Gemini, ChatGPT, and DeepSeek, noting that while DeepSeek had superior search capabilities for niche topics, server reliability issues led him to switch to other platforms. The conversation ended with DeanPretorius attempting to demonstrate something that wasn't working properly at the time.

\subsubsection{Physics Simulations Research Discussion}

The team discussed DeanPretorius's work on physics simulations in Google Colab, focusing on friction cones and Riemannian manifolds. Hozefa suggested exploring how multi-contact scenarios work with individual cones and manifolds for each contact point. DeanPretorius agreed to research the application of Christmanian manifolds in friction space and identify research gaps. The team agreed to review DeanPretorius's thesis and paper, with Fares requesting the documents be sent by email for further analysis before the next meeting.

 
 